% GENERATED FILE. Edit canonical lessons, then run npm run generate:books.
% canonical-source-hash: fnv1a64:b3fb245bdc668c2d
% canonical-lessons: UR-C477-qatra, UR-C477-gholna, UR-C477-kaipsul, UR-C477-puriya, UR-C477-phankna

\chapter{Taking Medicine}
\label{ch:ur-taking-medicine}
\hlchaptermodality{477}

\begin{chapteropening}
\textbf{By the end of this chapter:} \emph{I can say how to take a medicine.}

Complete the last lesson of chapter 477: I can say how to take a medicine.
\end{chapteropening}

\section[\texorpdfstring{\ur{قطرہ} (qatrā)}{قطرہ (qatrā)}]{\ur{قطرہ} (qatrā) — a droplet; a drop (of medicine)}

\label{lesson:UR-C477-qatra}



\begin{quote}
Before the new one: say the Urdu for a traditional physician, then the Urdu for to have (someone) sent for.
\end{quote}

\subsection*{You'll want to know: \ur{قطرہ}}
\textbf{\ur{قطرہ}} — \emph{qatrā} — \textquotedblleft{}a droplet; a drop (of medicine)\textquotedblright{}.

\begin{quote}
\ur{دوا} \ur{کے} \ur{دو} \ur{قطرے} \ur{پانی} \ur{میں} \ur{ملائیے۔}
\end{quote}

\begin{quote}
\emph{davā ke do qatre pānī meṅ milāiye.}
\end{quote}

\begin{quote}
Mix two drops of the medicine into some water.
\end{quote}

Masculine.

\begin{cousinweb}
\textbf{\ur{قطرہ}} (\emph{qatrā}) is usually traced, through Persian, to Arabic \emph{qaṭra}, \textquotedblleft{}a drop\textquotedblright{}.
\end{cousinweb}

\begin{grammarlens}[title={What you've built}]
One more word for describing what you lived and felt.
\end{grammarlens}

\subsection*{Your turn}
\begin{itemize}
\raggedright
  \item \emph{Say it:} \emph{qatrā}
  \item \emph{Say it:} \emph{qatrā}, once more
  \item \emph{Say it:} \emph{davā ke do qatre pānī meṅ milāiye.} — Mix two drops of the medicine into some water.
  \item \emph{Recall:} say the Urdu for a traditional physician, then the Urdu for to have (someone) sent for, then say \emph{qatrā} again
\end{itemize}

\begin{tcolorbox}[breakable,colback=teal!4,colframe=teal!35!black,title={Before you move on}]
What does \ur{قطرہ} mean? (\textquotedblleft{}A droplet; a drop (of medicine)\textquotedblright{}.) Say it once more.
\end{tcolorbox}
\section[\texorpdfstring{\ur{گھولنا} (gholnā)}{گھولنا (gholnā)}]{\ur{گھولنا} (gholnā) — to dissolve; to stir in}

\label{lesson:UR-C477-gholna}



\begin{quote}
Before the new one: say the Urdu for to have (someone) sent for, then the Urdu for a droplet.
\end{quote}

\subsection*{You'll want to know: \ur{گھولنا}}
\textbf{\ur{گھولنا}} — \emph{gholnā} — \textquotedblleft{}to dissolve; to stir in\textquotedblright{}.

\begin{quote}
\ur{امی} \ur{نے} \ur{دودھ} \ur{میں} \ur{چینی} \ur{گھولی۔}
\end{quote}

\begin{quote}
\emph{ammī ne dūdh meṅ chīnī gholī.}
\end{quote}

\begin{quote}
Mum stirred sugar into the milk.
\end{quote}

Past, with \ur{نے} (he or she …): \ur{اس} \ur{نے} \ur{پانی} \ur{میں} \ur{نمک} \ur{گھولا}, \emph{us ne pānī meṅ namak gholā}.

\begin{cousinweb}
\textbf{\ur{گھولنا}} (\emph{gholnā}) is inherited from Sanskrit \emph{ghōlayati}, \textquotedblleft{}mixes, stirs together\textquotedblright{} (CDIAL 4526).
\end{cousinweb}

\begin{grammarlens}[title={What you've built}]
One more word for describing what you lived and felt.
\end{grammarlens}

\subsection*{Your turn}
\begin{itemize}
\raggedright
  \item \emph{Say it:} \emph{gholnā}
  \item \emph{Say it:} \emph{gholnā}, once more
  \item \emph{Say it:} \emph{ammī ne dūdh meṅ chīnī gholī.} — Mum stirred sugar into the milk.
  \item \emph{Recall:} say the Urdu for to have (someone) sent for, then the Urdu for a droplet, then say \emph{gholnā} again
\end{itemize}

\begin{tcolorbox}[breakable,colback=teal!4,colframe=teal!35!black,title={Before you move on}]
What does \ur{گھولنا} mean? (\textquotedblleft{}To dissolve; to stir in\textquotedblright{}.) Say it once more.
\end{tcolorbox}
\section[\texorpdfstring{\ur{کیپسول} (kaipsūl)}{کیپسول (kaipsūl)}]{\ur{کیپسول} (kaipsūl) — a capsule}

\label{lesson:UR-C477-kaipsul}



\begin{quote}
Before the new one: say the Urdu for a droplet, then the Urdu for to dissolve.
\end{quote}

\subsection*{You'll want to know: \ur{کیپسول}}
\textbf{\ur{کیپسول}} — \emph{kaipsūl} — \textquotedblleft{}a capsule\textquotedblright{}.

\begin{quote}
\ur{صبح} \ur{ایک} \ur{کیپسول} \ur{پانی} \ur{کے} \ur{ساتھ} \ur{لیجیے۔}
\end{quote}

\begin{quote}
\emph{subah ek kaipsūl pānī ke sāth lījiye.}
\end{quote}

\begin{quote}
Take one capsule with water in the morning.
\end{quote}

Masculine.

\begin{cousinweb}
\textbf{\ur{کیپسول}} (\emph{kaipsūl}) is borrowed from English \emph{capsule}.
\end{cousinweb}

\begin{grammarlens}[title={What you've built}]
One more word for describing what you lived and felt.
\end{grammarlens}

\subsection*{Your turn}
\begin{itemize}
\raggedright
  \item \emph{Say it:} \emph{kaipsūl}
  \item \emph{Say it:} \emph{kaipsūl}, once more
  \item \emph{Say it:} \emph{subah ek kaipsūl pānī ke sāth lījiye.} — Take one capsule with water in the morning.
  \item \emph{Recall:} say the Urdu for a droplet, then the Urdu for to dissolve, then say \emph{kaipsūl} again
\end{itemize}

\begin{tcolorbox}[breakable,colback=teal!4,colframe=teal!35!black,title={Before you move on}]
What does \ur{کیپسول} mean? (\textquotedblleft{}A capsule\textquotedblright{}.) Say it once more.
\end{tcolorbox}
\section[\texorpdfstring{\ur{پڑیا} (puṛiyā)}{پڑیا (puṛiyā)}]{\ur{پڑیا} (puṛiyā) — a paper packet (of powder)}

\label{lesson:UR-C477-puriya}



\begin{quote}
Before the new one: say the Urdu for to dissolve, then the Urdu for a capsule.
\end{quote}

\subsection*{You'll want to know: \ur{پڑیا}}
\textbf{\ur{پڑیا}} — \emph{puṛiyā} — \textquotedblleft{}a paper packet (of powder)\textquotedblright{}.

\begin{quote}
\ur{حکیم} \ur{نے} \ur{دوا} \ur{کی} \ur{تین} \ur{پڑیاں} \ur{دیں۔}
\end{quote}

\begin{quote}
\emph{hakīm ne davā kī tīn puṛiyāṅ dīṅ.}
\end{quote}

\begin{quote}
The physician gave three packets of powder.
\end{quote}

Feminine.

\begin{cousinweb}
\textbf{\ur{پڑیا}} (\emph{puṛiyā}) is inherited: CDIAL lists the Hindi cousin \emph{puṛiyā}, \textquotedblleft{}anything wrapped in paper or leaves\textquotedblright{}, under Sanskrit \emph{puṭa}, \textquotedblleft{}a fold, a small receptacle\textquotedblright{} (CDIAL 8253).
\end{cousinweb}

\begin{grammarlens}[title={What you've built}]
One more word for describing what you lived and felt.
\end{grammarlens}

\subsection*{Your turn}
\begin{itemize}
\raggedright
  \item \emph{Say it:} \emph{puṛiyā}
  \item \emph{Say it:} \emph{puṛiyā}, once more
  \item \emph{Say it:} \emph{hakīm ne davā kī tīn puṛiyāṅ dīṅ.} — The physician gave three packets of powder.
  \item \emph{Recall:} say the Urdu for to dissolve, then the Urdu for a capsule, then say \emph{puṛiyā} again
\end{itemize}

\begin{tcolorbox}[breakable,colback=teal!4,colframe=teal!35!black,title={Before you move on}]
What does \ur{پڑیا} mean? (\textquotedblleft{}A paper packet (of powder)\textquotedblright{}.) Say it once more.
\end{tcolorbox}
\section[\texorpdfstring{\ur{پھانکنا} (phāṅknā)}{پھانکنا (phāṅknā)}]{\ur{پھانکنا} (phāṅknā) — to toss into the mouth (a dry powder)}

\label{lesson:UR-C477-phankna}



\begin{quote}
Before the new one: say the Urdu for a capsule, then the Urdu for a paper packet.
\end{quote}

\subsection*{You'll want to know: \ur{پھانکنا}}
\textbf{\ur{پھانکنا}} — \emph{phāṅknā} — \textquotedblleft{}to toss into the mouth (a dry powder)\textquotedblright{}.

\begin{quote}
\ur{اس} \ur{نے} \ur{پڑیا} \ur{کی} \ur{دوا} \ur{پھانکی۔}
\end{quote}

\begin{quote}
\emph{us ne puṛiyā kī davā phāṅkī.}
\end{quote}

\begin{quote}
He tossed the powder from the packet into his mouth.
\end{quote}

Past, with \ur{نے} (he or she …): \ur{اس} \ur{نے} \ur{چنے} \ur{پھانکے}, \emph{us ne chane phāṅke}.

\begin{cousinweb}
\textbf{\ur{پھانکنا}} (\emph{phāṅknā}) is inherited: CDIAL lists the Hindi cousin \emph{phā̃knā} under \emph{phakk}, \textquotedblleft{}toss food into the mouth\textquotedblright{}, a form it reconstructs (CDIAL 9034).
\end{cousinweb}

\begin{grammarlens}[title={What you've built}]
One more word for describing what you lived and felt.
\end{grammarlens}

\subsection*{Your turn}
\begin{itemize}
\raggedright
  \item \emph{Say it:} \emph{phāṅknā}
  \item \emph{Say it:} \emph{phāṅknā}, once more
  \item \emph{Say it:} \emph{us ne puṛiyā kī davā phāṅkī.} — He tossed the powder from the packet into his mouth.
  \item \emph{Recall:} say the Urdu for a capsule, then the Urdu for a paper packet, then say \emph{phāṅknā} again
\end{itemize}

\begin{tcolorbox}[breakable,colback=teal!4,colframe=teal!35!black,title={Before you move on}]
What does \ur{پھانکنا} mean? (\textquotedblleft{}To toss into the mouth (a dry powder)\textquotedblright{}.) Say it once more.
\end{tcolorbox}
